\documentclass[11pt,a4paper]{article}

\usepackage[utf8]{inputenc}
\usepackage[T1]{fontenc}
\usepackage{lmodern}
\usepackage[margin=2.2cm]{geometry}
\usepackage{amsmath,amssymb}
\usepackage{booktabs}
\usepackage{array}
\usepackage{longtable}
\usepackage{enumitem}
\usepackage{xcolor}
\usepackage{hyperref}
\usepackage[numbers,sort&compress]{natbib}
\usepackage{microtype}
\setlength{\emergencystretch}{3em}
\usepackage{fancyhdr}

\definecolor{ceresgreen}{HTML}{1B5E20}
\definecolor{ceresaccent}{HTML}{2E7D32}
\definecolor{ceresmild}{HTML}{EEF5EE}

\hypersetup{
  colorlinks=true,
  linkcolor=ceresgreen,
  citecolor=ceresaccent,
  urlcolor=ceresaccent,
  pdftitle={Ceres Agronomic Physics & 3-Year Crop Rotation Engine: Scientific Methodology},
  pdfauthor={Ceres Project (NASA Space Apps Challenge 2026)}
}

\pagestyle{fancy}
\fancyhf{}
\lhead{\small\textsc{Ceres} --- Agronomic Modeling \& Rotation Methodology}
\rhead{\small NASA Space Apps Challenge 2026}
\cfoot{\small\thepage}
\renewcommand{\headrulewidth}{0.4pt}

\title{%
  \vspace{-1.0cm}
  {\large\color{ceresaccent}\textsc{Ceres Decision-Support Platform}}\\[0.35em]
  \textbf{Scientific Methodology, Physical Equations, and Uncertainty Boundaries of the 3-Year Crop Rotation Analytics Engine}
}
\author{Ceres Engineering \& Agronomic Modeling Documentation}
\date{Version 1.1 --- November 2026}

\begin{document}
\maketitle

\begin{abstract}
\textbf{Ceres} is an open-source agronomic decision-support platform developed for the NASA Space Apps Challenge 2026 to help farmers evaluate 3-year crop rotations worldwide. Its core architectural invariant is \emph{deterministic physics first, generative language models strictly for explanation}: every hydrologic, soil carbon, nitrogen, erosion, compaction, and economic output is computed deterministically in Python before any natural-language summary is generated. Wherever possible the equations come from published methods (FAO-56, FAO-33, RUSLE, H{\'e}nin--Dupuis, Bolinder, IPCC); every parameter without a direct literature source is listed explicitly as a calibration parameter (Table~\ref{tab:calibration}), and heuristic components are labelled as such. This technical note documents the mathematical formulations implemented in \texttt{analytics-api} and the upstream Earth-observation ingestion pipeline in \texttt{web-api}, records parameter sources and unit normalizations, and states physical and data limitations alongside a farmer-facing field verification protocol. Version~1.1 replaces the half-year water balance with a monthly crop-calendar balance, corrects the FAO-33 yield response, nitrogen-fixation and fertilizer-recovery terms, and documents regional data sources for the Turin and Istanbul demonstration sites.
\end{abstract}

\tableofcontents
\vspace{0.8em}
\hrule
\vspace{0.8em}

\section{Architecture \& Data Layering}
\label{sec:architecture}

\subsection{Separation of Physics and Natural-Language Explanation}
The Ceres analytics pipeline enforces strict separation between deterministic simulation and generative language synthesis across two services:
\begin{enumerate}[leftmargin=1.5em]
  \item \textbf{Upstream Earth-Observation Preparation (\texttt{web-api})}: Aggregates parcel geometry, multi-source NASA and European environmental snapshots, soil profile grids, national market price indices, regional yield baselines, and paired ESP32 field probe measurements via \texttt{build\_ingestion\_payload} before calling the analytics endpoint.
  \item \textbf{Deterministic Physics \& Optimization (\texttt{analytics-api})}: Normalizes heterogeneous inputs in \texttt{FieldIngestionParser}, resolves a strictly ordered provenance hierarchy, runs coupled annual soil--water--crop simulations in \texttt{SoilPhysicsEngine}, searches all admissible 3-year crop sequences in \texttt{PlanOptimizerEngine}, and returns five non-redundant strategic champions.
  \item \textbf{Downstream Advisory Narration (\texttt{GeminiAdvisorEngine})}: Consumes only the completed JSON outputs and climate context to translate numerical trade-offs and physical alerts into farmer-readable guidance. Large-language-model calls have \emph{zero} write access to yields, water deficits, nitrogen balances, or strategy rankings; whenever the external LLM service is unavailable or times out, a deterministic Python fallback formatter generates structured rationales directly from the simulated metrics.
\end{enumerate}

\subsection{Priority Hierarchy and Provenance Resolution}
Field-scale observations frequently conflict with coarse-resolution satellite grids or regional cadastral defaults. Ceres resolves every agronomic parameter according to an explicit three-tier override hierarchy implemented in \texttt{FieldIngestionParser.parse\_and\_validate}:

\begin{table}[h!]
  \centering
  \small
  \caption{Three-tier data precedence rules implemented in \texttt{FieldIngestionParser}.}
  \label{tab:data-layering}
  \begin{tabular}{@{}p{2.2cm}p{4.6cm}p{7.4cm}@{}}
    \toprule
    \textbf{Tier} & \textbf{Source Class} & \textbf{Overridden Parameters \& Handling} \\
    \midrule
    Tier 1 (Highest) & User / Farmer Overrides \& Parcel Covenants & Soil texture class, measured SOC~(\%), measured $\mathrm{pH}$, slope~(\%), rain-gauge annual precipitation, penetrometer cone index~($\mathrm{PSI}$), annual irrigation water quota~($\mathrm{m^3/yr}$), cropping whitelist/blacklist, greenhouse infrastructure, low-spot flood risk, recent 5-year crop history. \\
    \addlinespace
    Tier 2 & Dated In-Situ ESP32 Telemetry & Chronologically sorted sensor logs (soil volumetric moisture, soil temperature, electrical conductivity, optional penetrometer $\mathrm{PSI}$, and paired canopy ambient air temperature \& relative humidity). The latest valid observation seeds initial root-zone moisture and microclimate adjustment. \\
    \addlinespace
    Tier 3 (Baseline) & Gridded Earth Observation \& Public Statistics & NASA POWER agroclimatology (1991--2020 monthly normals and daily series)~\cite{nasapower2024}, NASA-USDA / SMAP root-zone moisture, ISRIC SoilGrids 2.0 stratified horizons~\cite{poggio2021soilgrids}, Copernicus/SRTM terrain slope, FAOSTAT/Eurostat/World Bank regional yields and prices~\cite{faostat2024}. \\
    \bottomrule
  \end{tabular}
\end{table}

\subsection{Coordinate, Depth-Profile, and Unit Normalization}
Before running physics kernels, ingestion normalizes units and spatial encodings:
\begin{itemize}[leftmargin=1.5em]
  \item \textbf{Geometry and Area}: Parcel GeoJSON rings (including cadastral polygons transformed from EPSG:6706 \texttt{(lat, lon)} to standard \texttt{[lon, lat]}) yield polygon centroid latitude/longitude and geodesic parcel area in hectares ($\mathrm{ha}$; farmer-entered acre inputs are converted via $1~\mathrm{ac} = 0.404686~\mathrm{ha}$).
  \item \textbf{Stratified Soil Horizons}: ISRIC SoilGrids 2.0 supplies properties over depth intervals ($0$--$5$, $5$--$15$, $15$--$30$, $30$--$60$, $60$--$100~\mathrm{cm}$)~\cite{poggio2021soilgrids}. Ceres normalizes SoilGrids scaled integers into physical units: soil $\mathrm{pH} = \mathrm{phh2o} / 10$, bulk density $\rho_b = \mathrm{bdod} / 100~\mathrm{g\,cm^{-3}}$, soil organic carbon $\mathrm{SOC}_{\mathrm{\%}} = (\mathrm{soc}_{\mathrm{dg/kg}} / 10) / 10 = \mathrm{g\,kg^{-1}} / 10$, and USDA texture mass fractions ($\mathrm{clay}, \mathrm{sand}, \mathrm{silt}$ in~$\%$) from $\mathrm{g\,kg^{-1}} / 10$.
  \item \textbf{Hydrologic \& Economic Normalization}: Irrigation availability scales from legal parcel quota~$Q_{\mathrm{yr}}$~($\mathrm{m^3\,yr^{-1}}$) and parcel area~$A_{\mathrm{ha}}$ to equivalent water depth $Q_{\mathrm{mm}} = Q_{\mathrm{yr}} / (10 \, A_{\mathrm{ha}})$. Nutrient budgets track strictly $\mathrm{kg~N~ha^{-1}}$ and crop prices track $\mathrm{USD~t^{-1}}$ ($\mathrm{metric~tonne}^{-1}$).
  \item \textbf{Unmeasured State Provenance}: When local penetrometer readings are absent, baseline compaction is labeled explicitly in provenance metadata as \emph{``Unmeasured (modeled from machinery passes \& soil texture; verify with spade/penetrometer)''} so modeled inferences are never misrepresented as sensor measurements.
\end{itemize}

\section{Agroclimatology \& Field Hydrology}
\label{sec:hydrology}

All water-balance quantities are computed on a \emph{monthly} time step over the months each crop actually occupies, using a per-parcel crop calendar (Section~\ref{sec:calendar}). This replaces the fixed warm/cool half-year windows of Version~1.0, which assigned autumn-sown crops (winter wheat, barley, rapeseed) and perennials (alfalfa) to a single half-year and therefore omitted their spring peak water use.

\subsection{Monthly Reference Evapotranspiration \texorpdfstring{$ET_0$}{ET0}}
\label{sec:et0}
\paragraph{Primary method: monthly FAO-56 Penman--Monteith.}
NASA POWER 1991--2020 monthly climatology provides every variable required by the FAO-56 Penman--Monteith equation~\cite{allen1998fao56,nasapower2024}: mean daily maximum and minimum air temperature at 2~m (\texttt{T2M\_MAX}, \texttt{T2M\_MIN}), relative humidity (\texttt{RH2M}), dew-point temperature (\texttt{T2MDEW}), wind speed at 2~m (\texttt{WS2M}) and all-sky shortwave radiation (\texttt{ALLSKY\_SFC\_SW\_DWN}). Monthly $ET_{0,m}$ ($\mathrm{mm\,day^{-1}}$) is therefore computed with FAO-56 Eq.~6:
\begin{equation}
  ET_0 = \frac{0.408\,\Delta\,(R_n - G) + \gamma\,\dfrac{900}{T_{\mathrm{mean}} + 273}\, u_2 \,(e_s - e_a)}{\Delta + \gamma\,(1 + 0.34\, u_2)},
\end{equation}
with atmospheric pressure $P_{\mathrm{atm}} = 101.3 \bigl(\frac{293 - 0.0065 z}{293}\bigr)^{5.26}$~($\mathrm{kPa}$, Eq.~7), psychrometric constant $\gamma = 0.665\times 10^{-3} P_{\mathrm{atm}}$ (Eq.~8), $e_s = \tfrac{1}{2}[e^{\circ}(T_{\max}) + e^{\circ}(T_{\min})]$ (Eq.~12), and actual vapour pressure from dew point, $e_a = e^{\circ}(T_{\mathrm{dew}})$ (Eq.~14), falling back to $e_a = e_s\,\mathrm{RH}_{\mathrm{mean}}/100$ (Eq.~19) only when dew point is unavailable. Net radiation is $R_n = (1-\alpha)R_s - R_{nl}$ with $\alpha = 0.23$, and net longwave radiation uses the \emph{actual} relative shortwave radiation (FAO-56 Eq.~39):
\begin{equation}
  R_{nl} = \sigma \left[\frac{T_{\max,K}^4 + T_{\min,K}^4}{2}\right]\bigl(0.34 - 0.14\sqrt{e_a}\bigr)\left(1.35\,\frac{R_s}{R_{so}} - 0.35\right), \qquad R_{so} = (0.75 + 2\times10^{-5} z)\,R_a,
\end{equation}
with $R_s/R_{so}$ limited to $\le 1.0$. Monthly soil heat flux is not neglected: $G_m = 0.07\,(T_{\mathrm{mean},m+1} - T_{\mathrm{mean},m-1})$ (FAO-56 Eq.~43; climatological months are cyclic, so the following month is always known). Extraterrestrial radiation $R_a$ on the mid-month day of year $J_m$ follows FAO-56 Eqs.~21--25:
\begin{align}
  d_r &= 1 + 0.033 \cos\!\left(\frac{2\pi}{365} J_m\right), &
  \delta &= 0.409 \sin\!\left(\frac{2\pi}{365} J_m - 1.39\right), \\
  \omega_s &= \arccos\!\bigl(\operatorname{clamp}(-\tan\phi \tan\delta,\, -1,\, 1)\bigr), &
  R_a &= \frac{24(60)}{\pi} G_{sc}\, d_r \Bigl(\omega_s \sin\phi\sin\delta + \cos\phi\cos\delta\sin\omega_s\Bigr),
\end{align}
where $G_{sc} = 0.0820~\mathrm{MJ\,m^{-2}\,min^{-1}}$ and $\phi$ is latitude (rad). The same kernel is used on a daily step (\texttt{calculate\_\allowbreak fao56\_\allowbreak penman\_\allowbreak monteith\_et0}) when daily series are supplied, with $G \approx 0$ (FAO-56 Eq.~42).

\paragraph{Fallback: Hargreaves--Samani.}
Only when humidity, wind or radiation are missing is $ET_{0,m}$ estimated with FAO-56 Eq.~52~\cite{hargreaves1985reference,allen1998fao56}:
\begin{equation}
  ET_{0,m} = 0.0023 \, \bigl(T_{\mathrm{mean},m} + 17.8\bigr) \sqrt{\max\!\bigl(0.1,\, T_{\max,m} - T_{\min,m}\bigr)} \, \bigl(0.408\, R_a\bigr).
\end{equation}
The provenance metadata records which method was used. FAO-56 notes that Eq.~52 tends to over-predict under high relative humidity and under-predict under high wind; both conditions occur at the demonstration sites (humid Po plain near Turin; windy Bosphorus coast near Istanbul), which is why Penman--Monteith is the primary method.

\subsection{Crop Calendar and Monthly Crop Coefficients}
\label{sec:calendar}
\paragraph{Calendar.} Sowing and harvest months for each crop are resolved with the same three-tier precedence as other parameters: farmer-reported dates (Tier~1), otherwise the GGCMI Phase~3 global crop calendar at $0.5^{\circ}$ (planting and maturity day per crop, rainfed and irrigated separately)~\cite{jagermeyr2021calendar}, otherwise catalog defaults. A crop occupies month $m$ with fraction $f_m \in [0,1]$ (partial first and last months).

\paragraph{Single crop coefficient curve.} Crop evapotranspiration uses the FAO-56 \emph{single} crop coefficient $K_c$ (Table~12), not the basal coefficient $K_{cb}$, because the lumped monthly balance does not model soil evaporation separately. Stage lengths $L_{\mathrm{ini}}, L_{\mathrm{dev}}, L_{\mathrm{mid}}, L_{\mathrm{late}}$ (days) come from FAO-56 Table~11 scaled to the resolved season length. Following FAO-56, $K_c$ is constant at $K_{c,\mathrm{ini}}$ during the initial stage, rises linearly to $K_{c,\mathrm{mid}}$ during development, is constant during mid-season and falls linearly to $K_{c,\mathrm{end}}$ during late season. The monthly coefficient $K_{c,m}$ is the mean of this daily curve over the crop days in month~$m$. $K_{c,\mathrm{mid}}$ is adjusted for climate with FAO-56 Eq.~62,
\begin{equation}
  K_{c,\mathrm{mid}} = K_{c,\mathrm{mid}}^{\mathrm{Tab}} + \bigl[0.04\,(u_2 - 2) - 0.004\,(\mathrm{RH}_{\min} - 45)\bigr]\left(\frac{h}{3}\right)^{0.3},
\end{equation}
where $h$ is mean crop height (m) and $\mathrm{RH}_{\min} = 100\, e^{\circ}(T_{\mathrm{dew}})/e^{\circ}(T_{\max})$ during mid-season. Perennial forage (alfalfa) uses the FAO-56 averaged cutting-cycle coefficient for every month of the stand. Monthly crop evapotranspiration is
\begin{equation}
  ET_{c,m} = K_{c,m}\, ET_{0,m}\, N_m\, f_m ,
\end{equation}
with $N_m$ the number of days in month $m$.

\subsection{Effective Rainfall}
Monthly effective rainfall uses the USDA-SCS formula as implemented in FAO CROPWAT~\cite{smith1992cropwat}, applied to each \emph{monthly} total $P_m$ (mm):
\begin{equation}
  P_{\mathrm{eff},m} =
  \begin{cases}
    P_m\,\dfrac{125 - 0.2\,P_m}{125}, & P_m \le 250~\mathrm{mm}, \\[6pt]
    125 + 0.1\,P_m, & P_m > 250~\mathrm{mm}.
  \end{cases}
\end{equation}
Version~1.0 used $\max(0, 0.75\,P_{\mathrm{seas}} - 10)$ on half-year totals. That expression is not the USDA-SCS method, and its offset is meaningful only for monthly totals.

\subsection{Monthly Root-Zone Water Balance}
Total and readily available water follow FAO-56 Eqs.~82--83:
\begin{equation}
  \mathrm{TAW} = 1000\,(\theta_{\mathrm{fc}} - \theta_{\mathrm{wp}})\, Z_r, \qquad
  \mathrm{RAW} = p\,\mathrm{TAW}, \qquad p = p_{\mathrm{Tab}} + 0.04\,(5 - ET_c),
\end{equation}
with rooting depth $Z_r$ and depletion fraction $p_{\mathrm{Tab}}$ from FAO-56 Table~22 ($p$ bounded to $[0.1, 0.8]$; $ET_c$ in $\mathrm{mm\,day^{-1}}$). Field capacity and wilting point are estimated from local or SoilGrids sand, clay and organic matter with the Saxton--Rawls pedotransfer functions~\cite{saxton2006soil}, falling back to FAO-56 Table~19 texture-class midpoints.

Root-zone storage $S$ (mm, $0 \le S \le \mathrm{TAW}$) is carried month by month and across crop years. With start-of-month depletion $D_{r,m} = \mathrm{TAW} - S_{m-1}$, the FAO-56 water-stress coefficient (Eq.~84) is applied to the \emph{root-zone depletion state}, as FAO-56 defines it:
\begin{equation}
  K_{s,m} =
  \begin{cases}
    1, & D_{r,m} \le \mathrm{RAW}, \\[4pt]
    \dfrac{\mathrm{TAW} - D_{r,m}}{(1 - p)\,\mathrm{TAW}}, & D_{r,m} > \mathrm{RAW},
  \end{cases}
  \qquad
  ET_{a,m} = \min\!\bigl(K_{s,m}\, ET_{c,m},\; S_{m-1} + P_{\mathrm{eff},m} + I_{\mathrm{net},m}\bigr),
\end{equation}
\begin{equation}
  S_m = \min\!\bigl(\mathrm{TAW},\; S_{m-1} + P_{\mathrm{eff},m} + I_{\mathrm{net},m} - ET_{a,m}\bigr),
\end{equation}
water above $\mathrm{TAW}$ being treated as drainage. In months without a crop, the parcel is treated as bare or residue-covered soil with $K_{c} = K_{c,\mathrm{ini}}$ of the following crop, which represents mainly soil evaporation in FAO-56.

\paragraph{Initial state.} Year~1 starts from the state obtained by running the balance through the 12 months preceding sowing with the farmer-reported Year~0 crop. Without crop history, $S_0 = 1000\,\bigl(\min(\theta_0, \theta_{\mathrm{fc}}) - \theta_{\mathrm{wp}}\bigr)^+ Z_r$, with $\theta_0$ the latest probe reading taken within 30 days of sowing, or else the NASA SMAP Level-4 root-zone soil moisture climatology for the month before sowing. The ad-hoc factor $0.4$ of Version~1.0 is removed.

\paragraph{Irrigation.} When irrigation is enabled, net irrigation in month $m$ refills the root zone to the readily-available threshold, $I_{\mathrm{net},m} = \max\bigl(0,\, ET_{c,m} - P_{\mathrm{eff},m} - (S_{m-1} - (\mathrm{TAW} - \mathrm{RAW}))\bigr)$, limited by the remaining annual quota. Gross withdrawal is $I_{\mathrm{gross},m} = I_{\mathrm{net},m}/e_a$. The field application efficiency $e_a$ depends on the irrigation method: $0.60$ for surface (border, furrow, basin), $0.75$ for sprinkler and $0.90$ for drip~\cite{fao_irrigation_efficiency}. The annual cap is $\sum_m I_{\mathrm{gross},m} \le Q_{\mathrm{mm}} = Q_{\mathrm{yr}}/(10\,A_{\mathrm{ha}})$, shared by every crop grown in that calendar year. Gravity-fed surface irrigation from canals is common in the Turin area, so declaring the system type matters: at $e_a = 0.60$ the same quota delivers about 30\% less water to the root zone than the single value $0.85$ used in Version~1.0.

\subsection{Seasonal Yield Response to Water (FAO-33)}
Water-limited yield follows the FAO-33 seasonal relationship between relative yield and relative evapotranspiration~\cite{doorenbos1979fao33}, evaluated with the seasonal sums of the monthly balance:
\begin{equation}
  1 - \frac{Y_a}{Y_m} = K_y \left(1 - \frac{\sum_m ET_{a,m}}{\sum_m ET_{c,m}}\right), \qquad
  \mu_{\mathrm{water}} = \max\!\left(0,\; 1 - K_y \left(1 - \frac{\sum_m ET_{a,m}}{\sum_m ET_{c,m}}\right)\right),
\end{equation}
with crop-specific seasonal $K_y$ from FAO-33. The linear relationship was derived for moderate deficits; when the relative evapotranspiration deficit exceeds $0.5$ the result is flagged as \emph{outside the calibrated range} instead of being silently clamped. Version~1.0 substituted the seasonal unmet demand for the daily depletion state in $K_s$ and then applied $K_y(1-K_s)$, which roughly doubled the predicted loss for moderate deficits; the arbitrary floor ($K_s \ge 0.30$) and cap ($\Delta Y \le 0.65$) that compensated for this are removed.

Observed daily climate extremes (\texttt{frost\_\allowbreak days\_\allowbreak count}, \texttt{heatwave\_\allowbreak days\_\allowbreak count}, \texttt{longest\_\allowbreak dry\_\allowbreak spell\_\allowbreak days}) and ESP32 canopy temperature/humidity further modulate cold-kill pruning and heat-stress yield penalties (Section~\ref{sec:optimization}). A seasonal balance cannot resolve a short drought at flowering; this limitation is stated in Section~\ref{sec:limitations}.

\section{Soil Physics: Carbon, Nitrogen, Erosion, and Compaction}
\label{sec:soil-physics}

\subsection{Soil Organic Carbon (H{\'e}nin--Dupuis / Mary--Gu{\'e}rif / Bolinder)}
Topsoil organic carbon evolution in the upper $0.30~\mathrm{m}$ arable horizon follows the classic first-order mass-balance differential equation of H{\'e}nin and Dupuis~\cite{henin1945essai}:
\begin{equation}
  \frac{\mathrm{d}C}{\mathrm{d}t} = k_1 A - k_2 C,
\end{equation}
where initial carbon stock $C_0$ ($\mathrm{t\,C\,ha^{-1}}$) is computed from topsoil bulk density $\rho_b$ ($\mathrm{g\,cm^{-3}}$) and organic carbon concentration $\mathrm{SOC}_0$ ($\%$):
\begin{equation}
  C_0 = 100 \cdot \rho_b \cdot 0.30 \cdot \mathrm{SOC}_0.
\end{equation}

The annual mineralization rate constant $k_2$ ($\mathrm{yr^{-1}}$) implements the Mary--Gu{\'e}rif temperature and clay-protection parameterization~\cite{mary1994soil}, scaled by tillage intensity factor $f_{\mathrm{till}} \in \{0.50, 0.70, 1.00\}$ (perennial/no-till, strip-till, conventional tillage):
\begin{equation}
  k_2 = f_{\mathrm{till}} \cdot \max\!\left(0.012,\; \frac{1.2 \max(6.0, T_{\mathrm{mean}}) - 6.0}{10\,\mathrm{Clay}_{\%} + 200}\right).
\end{equation}
Fresh biomass carbon inputs $A$ ($\mathrm{t\,C\,ha^{-1}\,yr^{-1}}$) follow Bolinder et al.~\cite{bolinder2007approach} Tier-2 allometric allocation using simulated harvested yield $Y$ ($\mathrm{t\,ha^{-1}}$), dry-matter fraction $f_{\mathrm{DM}}$ ($0.14$ for high-moisture horticultural crops such as tomatoes, bell peppers, and spinach; $0.86$ for grains, oilseeds, and dry forage), harvest index $\mathrm{HI}$, shoot-to-root ratio, and humification coefficient $k_1$:
\begin{equation}
  C_{\mathrm{prod}} = Y \cdot f_{\mathrm{DM}} \cdot 0.45, \qquad
  A = C_{\mathrm{stubble}} + C_{\mathrm{roots}} + C_{\mathrm{exudates}} + C_{\mathrm{cover}},
\end{equation}
integrated over an annual step ($\Delta C = k_1 A - k_2 C_t$, bounded to $[-3.5, +4.5]~\mathrm{t\,C\,ha^{-1}\,yr^{-1}}$) and converted back to topsoil $\mathrm{SOC}_\%$ after each cropping year.

\subsection{Biological Nitrogen Fixation and Seasonal N Budget}
Ceres computes nitrogen dynamics separately for leguminous and non-leguminous species:
\begin{itemize}[leftmargin=1.5em]
  \item \textbf{Legumes (Herridge \& Peoples Tier-2 BNF~\cite{herridge2008global})}: Total crop nitrogen uptake $N_{\mathrm{plant}} = 1.35 \, Y \, u_N$ (including non-harvested roots and nodules, where $u_N$ is harvested $\mathrm{kg~N~t^{-1}}$) yields biologically fixed nitrogen $N_{\mathrm{fix}} = N_{\mathrm{plant}} \cdot \mathrm{\%Ndfa}$. Net soil balance is $N_{\mathrm{fix}} - N_{\mathrm{harvest}} + N_{\mathrm{credit}}$, with synthetic fertilizer demand set to $0~\mathrm{kg~N~ha^{-1}}$.
  \item \textbf{Non-Legumes (Residual Carryover + Applied Synthetic N)}: Available carry-over combines any preceding legume nitrogen credit ($N_{\mathrm{leg,credit}}$), taproot cover-crop scavenged nitrogen ($N_{\mathrm{scav}}$), and $30\%$ of the mineralized topsoil nitrogen pool (capped at $45\%$ of crop uptake demand). Required synthetic nitrogen is then:
  \begin{equation}
    N_{\mathrm{synth}} = \max\!\bigl(0,\; N_{\mathrm{demand}} - N_{\mathrm{leg,used}} - N_{\mathrm{scav}} - N_{\mathrm{min,used}}\bigr),
  \end{equation}
  with annual net soil-N balance tracked as $\Delta N = \max\!\bigl(-80,\; N_{\mathrm{synth}} + N_{\mathrm{leg,used}} + N_{\mathrm{scav}} - N_{\mathrm{demand}}\bigr)$~($\mathrm{kg~N\,ha^{-1}}$).
\end{itemize}

\subsection{Topsoil Erosion via RUSLE}
Annual sheet and rill soil loss $E$ ($\mathrm{t\,ha^{-1}\,yr^{-1}}$) is computed from the Revised Universal Soil Loss Equation (RUSLE)~\cite{renard1997rusle}:
\begin{equation}
  E = R \cdot K \cdot LS \cdot C \cdot P,
\end{equation}
using four physically parameterized factors and a conservative support baseline ($P = 1.0$, unless specific structural practices are encoded):
\begin{enumerate}[leftmargin=1.5em]
  \item \textbf{Rainfall Erosivity $R$ ($\mathrm{MJ\,mm\,ha^{-1}\,h^{-1}\,yr^{-1}}$)} via Renard \& Freimund~\cite{renard1994using} from mean annual precipitation $P_{\mathrm{ann}}$ ($\mathrm{mm}$):
  \begin{equation}
    R =
    \begin{cases}
      0.04830 \, P_{\mathrm{ann}}^{1.610}, & P_{\mathrm{ann}} \le 850~\mathrm{mm}, \\
      587.8 - 1.219 \, P_{\mathrm{ann}} + 0.004105 \, P_{\mathrm{ann}}^{2}, & P_{\mathrm{ann}} > 850~\mathrm{mm}.
    \end{cases}
  \end{equation}
  \item \textbf{Soil Erodibility $K$ ($\mathrm{t\,ha\,h\,ha^{-1}\,MJ^{-1}\,mm^{-1}}$)} via the EPIC formulation of Williams~\cite{williams1995epic} using SoilGrids/local sand ($\mathrm{SAN}$), silt ($\mathrm{SIL}$), clay ($\mathrm{CLA}$), and organic carbon ($\mathrm{SOC}_\%$):
  \begin{equation}
    K = 0.1317 \cdot f_{\mathrm{csand}} \cdot f_{\mathrm{cl\text{-}si}} \cdot f_{\mathrm{orgc}} \cdot f_{\mathrm{hisand}},
  \end{equation}
  where $0.1317$ converts US customary units to SI.
  \item \textbf{Topographic Factor $LS$} via Moore \& Burch~\cite{moore1986physical}: slope angle $\theta = \arctan(\mathrm{slope}_{\%}/100)$ and effective slope flow path $\lambda = \max\bigl(22.13,\, \min(150.0,\, 0.5\sqrt{10^4 A_{\mathrm{ha}}})\bigr)$~($\mathrm{m}$) give:
  \begin{equation}
    LS = \left(\frac{\lambda}{22.13}\right)^{0.4} \left(\frac{\sin\theta}{0.0896}\right)^{1.3}.
  \end{equation}
  \item \textbf{Crop Cover-Management Factor $C$}: Assigned per crop profile ($C \in [0.02, 0.36]$, lowest for perennial alfalfa and winter cover/wheat, highest for wide-row tillage crops).
\end{enumerate}

\subsection{Soil Compaction Heuristic and Bio-Drill Relief}
Traffic compaction and taproot remediation are tracked in cone penetrometer equivalent units ($\mathrm{PSI}$; $1~\mathrm{MPa} \approx 145.038~\mathrm{PSI}$~\cite{asabe2019s313}). When field penetrometer measurements exist, the baseline starts at the measured $\mathrm{PSI}$; otherwise, a texture-calibrated prior is used and explicitly flagged as unmeasured. Each crop season increments $\mathrm{PSI}$ via an axle-pass curve $(8.0 + 4.0\min(10, N_{\mathrm{passes}}))(1 + \mathrm{Clay}_{\%}/100)$ and applies deep-taproot bio-drilling relief ($\Delta_{\mathrm{biodrill}} \in \{-65, -55, -25, -15\}~\mathrm{PSI}$ for Daikon radish, alfalfa, sunflowers, and canola), bounded to $[120, 380]~\mathrm{PSI}$.

\section{Crop Envelopes, Rotation Constraints, and Optimization}
\label{sec:optimization}

\subsection{Agronomic Crop Catalog and Biophysical Filtering}
\texttt{MASTER\_CROPS\_CATALOG} defines twelve representative field, forage, horticultural, and bio-drilling cover crops spanning seven botanical families (Poaceae, Fabaceae, Asteraceae, Brassicaceae, Solanaceae, Amaranthaceae). Before exhaustive sequence evaluation, \texttt{PlanOptimizerEngine} invokes \texttt{\_filter\_candidate\_crops} to prune inadmissible species using field constraints and climate observations:
\begin{enumerate}[leftmargin=1.5em]
  \item \textbf{Farm Covenants \& Infrastructure}: Blacklisted species are excluded; when at least three whitelisted crops are provided, search narrows to the whitelist. Crops requiring protected cultivation are filtered unless greenhouse infrastructure is flagged, and low-spot fields with elevated waterlogging exposure retain only crops with sufficient flood tolerance ($\ge 2.0$).
  \item \textbf{Thermal Envelope \& Frost/Heat Pruning}: Each crop carries an optimal thermal window $[T_{\mathrm{opt,min}}, T_{\mathrm{opt,max}}]$ and lethal frost threshold $T_{\mathrm{frost,kill}}$. Warm-season and cool-season crops are tested against observed seasonal temperatures, frost-day counts (\texttt{frost\_days\_count}), coldest night temperature, and heatwave counts (\texttt{heatwave\_days\_count}) from NASA POWER.
  \item \textbf{Edaphic $\mathrm{pH}$ Lockout Pruning}: Soil $\mathrm{pH}$ resolved from local tests or SoilGrids must lie within $\pm 1.6~\mathrm{pH}$ units of $[\mathrm{pH}_{\mathrm{opt,min}}, \mathrm{pH}_{\mathrm{opt,max}}]$; otherwise the crop is excluded from the feasible candidate set (falling back gracefully only if filtering leaves fewer than three candidates).
\end{enumerate}

\subsection{Rotational Break Rules and Sequential Carryover}
Every ordered 3-year triplet $(c_1, c_2, c_3)$ is simulated sequentially across Years $1 \to 2 \to 3$, carrying forward updated state $(C_t, N_{\mathrm{min},t}, \mathrm{PSI}_t, \mathrm{family}_{t-1})$ and preceding legume/scavenger credits. Agronomic sequence constraints enforce:
\begin{itemize}[leftmargin=1.5em]
  \item \textbf{Phytosanitary Break Intervals} (\texttt{rotational\_break\_needed\_years}): Solanaceae crops (\texttt{processing\_tomatoes}, \texttt{bell\_peppers}) requiring a $\ge 3$-year soil-borne disease break are limited to \emph{at most once} in any 3-year sequence and check any reported Year~$0$ preceding crop. Crops with $\ge 2$-year breaks prohibit consecutive same-family succession, and consecutive pure cover-crop seasons (\texttt{daikon\_radish\_cover}) are disallowed.
  \item \textbf{Continuous-Family Pathology Drag}: When adjacent crops share a botanical family without meeting ideal break length, expected yield suffers a $12\%$ continuous-family drag ($0.88$ multiplier), avoiding arbitrary fabricated dollar fines while reflecting agronomic yield depression.
\end{itemize}

\subsection{Composite Yield Attenuation and Economic Valuation}
For each year $t \in \{1, 2, 3\}$, baseline potential yield $Y_{0,t}$ (calibrated against FAOSTAT national production yields when available for the parcel's country) is scaled by multiplicative abiotic and biological stress factors:
\begin{equation}
  Y_t = Y_{0,t} \cdot \mu_{\mathrm{water}} \cdot \mu_{\mathrm{pH}} \cdot \mu_{\mathrm{thermal}} \cdot \mu_{\mathrm{compaction}} \cdot \mu_{\mathrm{family}},
\end{equation}
where:
\begin{itemize}[leftmargin=1.5em]
  \item $\mu_{\mathrm{water}} = \max(0.35,\, 1 - \Delta Y_{\mathrm{water}})$ derives directly from FAO-33 stress $K_y(1 - K_s)$ (Section~\ref{sec:hydrology}).
  \item $\mu_{\mathrm{pH}}$ applies a nutrient lockout attenuation of $0.14$ per $\mathrm{pH}$ unit outside $[\mathrm{pH}_{\mathrm{opt,min}}, \mathrm{pH}_{\mathrm{opt,max}}]$.
  \item $\mu_{\mathrm{thermal}}$ accounts for seasonal departure from $[T_{\mathrm{opt,min}}, T_{\mathrm{opt,max}}]$, acute NASA POWER heatwave exposure, and live ESP32 canopy ambient temperature departures when sensor telemetry is present.
  \item $\mu_{\mathrm{compaction}}$ attenuates root-limited yield when $\mathrm{PSI}_t$ exceeds the crop's critical penetrometer threshold.
\end{itemize}

Farm-gate revenue prices $Y_t$ using market benchmarks ingested from World Bank, FAOSTAT, and Eurostat indices~\cite{faostat2024} (falling back to catalog defaults only when external series are unavailable for that commodity). Synthetic fertilizer costs scale linearly with simulated nitrogen demand $N_{\mathrm{synth}}$, phosphorus removal, and potassium removal priced against World Bank urea, TSP, and KCl nutrient tariffs.

\subsection{Five-Strategy Combinatorial Selection}
\texttt{PlanOptimizerEngine.optimize\_all\_strategies} evaluates every admissible 3-year triplet and selects five non-redundant strategic champions suited to distinct farmer priorities:
\begin{enumerate}[leftmargin=1.5em]
  \item \textbf{Most Income (\texttt{most\_income})}: Maximizes cumulative 3-year net farm margin ($\mathrm{USD}$).
  \item \textbf{Best Soil Health (\texttt{best\_soil\_health})}: Maximizes weighted multi-year $\Delta\mathrm{SOC}$, compaction reduction ($\Delta\mathrm{PSI} < 0$), erosion suppression ($\mathrm{RUSLE}~E$), and biological nitrogen enrichment.
  \item \textbf{Least Effort (\texttt{least\_effort})}: Minimizes field labor hours, machinery passes, and irrigation pumping burden while preserving positive farm margin.
  \item \textbf{Highest ROI (\texttt{highest\_roi})}: Maximizes capital return ratio ($\text{Net Profit} / \text{Operating Cost}$), favoring low-input drought-hardy or legume-subsidized sequences.
  \item \textbf{Least Crop Loss Risk (\texttt{least\_crop\_loss\_risk})}: Minimizes composite multi-hazard risk index spanning unmet water deficit, observed frost/heat extreme vulnerability, continuous-family disease exposure, and soil compaction excess.
\end{enumerate}

\section{Scientific Assumptions, Calibration Parameters, and Field Verification}
\label{sec:limitations}

\subsection{Explicit Physical Modeling Assumptions}
Ceres' lumped monthly and annual equations simplify real agroecosystem dynamics in the following ways:
\begin{itemize}[leftmargin=1.5em]
  \item \textbf{Monthly Temporal Resolution}: Water is balanced monthly over each crop's calendar, not with daily soil-water profiles or dynamic phenology. A short drought during a critical flowering week can reduce yields more than the monthly balance indicates. The seasonal FAO-33 $K_y$ also ignores the higher sensitivity of flowering compared with vegetative stages.
  \item \textbf{Calendar Uncertainty}: Gridded crop calendars ($0.5^{\circ}$) represent typical regional practice; farmer-reported sowing and harvest dates should always override them.
  \item \textbf{Empirical RUSLE Parameterization}: Erosivity comes from a 1~km interpolated map (or, as a fallback, a US-fitted regression on annual rainfall). Slope length is estimated from parcel area under a square-plot geometry, and support practice $P = 1.0$. Contour plowing, grass waterways or terracing reduce real soil loss relative to these estimates.
  \item \textbf{Reference Evapotranspiration}: Monthly Penman--Monteith on $0.5^{\circ}$ NASA POWER grids smooths local wind and humidity; the Hargreaves--Samani fallback tends to over-predict under high humidity and under-predict under strong wind~\cite{allen1998fao56}.
  \item \textbf{Independent Multiplicative Stresses}: Yield stresses are multiplied as if independent, and are applied relative to a modelled reference field to avoid double counting with statistical yields.
  \item \textbf{Nitrogen}: Fertilizer recovery efficiency, residue nitrogen release and losses are represented with single seasonal coefficients. Nitrogen losses use IPCC Tier-1 factors, which are inventory averages, not field predictions.
  \item \textbf{Catalog Scope and Linear Economics}: The optimizer evaluates twelve benchmark species and assumes linear farm-gate pricing, independent of cooperative contracts, storage or seasonal price shocks.
\end{itemize}

\subsection{Calibration Parameters Without a Direct Literature Source}
Table~\ref{tab:calibration} lists every numeric parameter that is not taken directly from a cited method. These values are design choices or provisional calibrations. Results that depend on them are reported with their sensitivity (strategy ranking re-run with each parameter at its range limits), and they are priorities for local calibration.

\begin{table}[htbp]
  \centering
  \small
  \caption{Calibration parameters and their status.}
  \label{tab:calibration}
  \begin{tabular}{@{}p{4.2cm}p{2.3cm}p{2.4cm}p{6.2cm}@{}}
    \toprule
    \textbf{Parameter} & \textbf{Default} & \textbf{Used in} & \textbf{Status / evidence} \\
    \midrule
    Fertilizer N recovery $\mathrm{RE}_N$ & $0.40$ (range $0.3$--$0.6$) & $N_{\mathrm{synth}}$ & On-farm US maize $\approx 0.37$~\cite{cassman2002agroecosystems}; calibrate regionally \\
    Legume residue N release $f_{\mathrm{rel}}$ & catalog & $N_{\mathrm{credit}}$ & No source yet; use regional extension tables \\
    Cover-crop scavenged N $N_{\mathrm{scav}}$ & catalog & $N_{\mathrm{synth}}$ & No source yet \\
    Humification $k_1$ (above / below ground) & catalog & SOC & To be taken from AMG-type parameterizations \\
    Extra-root carbon fraction & $0.65$ & $A$ & Lower bound of the $65$--$100\%$ range in the literature \\
    Carbonate coefficient in $k_2$ & $0.3$ & $k_2$ & Verify against R{\'e}my \& Marin-Lafl{\`e}che (1974) \\
    Temperature floor in $f(T)$ & $T \ge 6^{\circ}\mathrm{C}$ & $k_2$, $N_{\mathrm{min}}$ & Design choice \\
    Family yield drag $\mu_{\mathrm{family}}$ & $0.88$ & $Y_t$ & No source; replace with crop-specific values \\
    pH attenuation & $0.14$ per unit & $\mu_{\mathrm{pH}}$ & No source \\
    pH lockout fallback & $\pm 1.6$ units & filtering & Only where no ECOCROP absolute range \\
    Flood-tolerance threshold & $\ge 2.0$ (1--5) & filtering & No source \\
    Compaction increments, reliefs, bounds & Section~\ref{sec:soil-physics} & $\mathrm{PSI}_t$ & Heuristic; indicative only \\
    Slope-length bounds & $22.13$--$150~\mathrm{m}$ & $LS$ & Geometric assumption \\
    Risk and soil-health index weights & see code & strategy selection & Design choice; reported with results \\
    \bottomrule
  \end{tabular}
\end{table}

\subsection{Metrics Intentionally Excluded or Labelled}
\begin{enumerate}[leftmargin=1.5em]
  \item \textbf{No Synthetic Sensor Claims}: Subsoil compaction without a penetrometer log is labelled as modeled from texture and machinery passes, never presented as a measurement.
  \item \textbf{No Monetary Fines}: Earlier prototypes assigned flat USD penalties per violated family break. These are replaced by hard phytosanitary exclusions and the yield drag $\mu_{\mathrm{family}}$, which is itself listed as a calibration parameter in Table~\ref{tab:calibration}.
  \item \textbf{LLM Output Containment}: Generative advisory text only interprets verified JSON simulation results; it cannot adjust numbers, invent yields, or override biophysical exclusions.
\end{enumerate}

\subsection{Demonstration Regions: Turin (Italy) and Istanbul (T{\"u}rkiye)}
For the live demonstration parcels, the following inputs have the largest effect on results and should be supplied or checked before the demo:
\begin{itemize}[leftmargin=1.5em]
  \item \textbf{Irrigation method}: many farms around Turin use gravity-fed surface irrigation from canals ($e_a = 0.60$). Declaring sprinkler or drip where it is not used would overstate delivered water by $25$--$50\%$.
  \item \textbf{Crop calendar}: autumn-sown cereals and rapeseed are common in both regions. Farmer-reported sowing and harvest months should override the gridded calendar.
  \item \textbf{Carbonate}: run an HCl effervescence test (or use LUCAS data for Italian parcels) wherever $\mathrm{pH} \ge 7.5$, because carbonate changes $k_2$ and therefore the soil-carbon trajectory.
  \item \textbf{Reference yields}: use ISTAT data for the Province of Turin and TurkStat data for Istanbul province rather than national averages, which mix very different growing regions.
  \item \textbf{Weather}: multi-year records from regional networks (e.g., the Piedmont agrometeorological network, the Turkish State Meteorological Service) can replace NASA POWER normals as Tier-1 data when they span at least 10 years.
  \item \textbf{Regulatory limits}: Nitrates Directive vulnerable-zone limits apply to parts of the Po plain.
\end{itemize}

\subsection{Human-in-the-Loop Verification Checklist}
Every strategy report includes a field-verification checklist (``Where Satellite Data Needs Local Eyes'') asking farmers and local agronomists to confirm parameters that coarse 250~m--0.5$^{\circ}$ remote sensing cannot resolve:
\begin{itemize}[leftmargin=1.5em]
  \item \textbf{Topsoil $\mathrm{pH}$, Carbonate \& Lab NPK}: Verify SoilGrids $\mathrm{pH}$ with composite 0--20~cm samples before high-value horticultural planting or liming; test for free carbonate where $\mathrm{pH} \ge 7.5$.
  \item \textbf{Tillage Pan Compaction}: Probe field diagonals with a spade or cone penetrometer after harvest; the modeled compaction trend is a heuristic.
  \item \textbf{Micro-Drainage \& Low Spots}: Confirm surface ponding, subsurface tile drainage capacity, and local wind-exposure microclimates during early establishment.
  \item \textbf{Irrigation Method \& Allocation}: Confirm the irrigation system and the seasonal allocation permit, since both set the water actually delivered.
  \item \textbf{Local Rotation History, Calendar \& Contracts}: Record crops of the past five years and usual sowing/harvest months, confirm buyer contracts and prices, and check regional fertilization limits before committing seed orders.
\end{itemize}


\bibliographystyle{plainnat}
\bibliography{references}

\end{document}
